\documentclass[../../main]{subfiles}

\begin{document}

\section{Relações de Ordem}
\label{sec:matematica:relacoes:relacoes-de-ordem}

\begin{defn}[Pré-Ordem]
  \label{def:matematica:relacoes:pre-ordem}

  Uma relação (R) sobre um conjunto (A) é uma pré-ordem quando é

  \begin{enumerate}
    \item reflexiva;
    \item transitiva.
  \end{enumerate}

\end{defn}

\begin{defn}[Ordem Parcial]
  \label{def:matematica:relacoes:ordem-parcial}

  Uma relação (R) sobre um conjunto (A) é uma ordem parcial quando
  é

  \begin{enumerate}
    \item reflexiva;
    \item antissimétrica;
    \item transitiva.
  \end{enumerate}

\end{defn}

\begin{defn}[Conjunto Parcialmente Ordenado]
  \label{def:matematica:relacoes:poset}

  Um conjunto parcialmente ordenado é um par

  [
    (A,R),
  ]

  onde (R) é uma ordem parcial sobre (A).

\end{defn}

\begin{defn}[Elementos Comparáveis]
  \label{def:matematica:relacoes:comparaveis}

  Dois elementos (a,b\in A) são comparáveis quando

  [
    aRb
  ]

  ou

  [
    bRa.
  ]

\end{defn}

\begin{defn}[Ordem Total]
  \label{def:matematica:relacoes:ordem-total}

  Uma ordem parcial (R) é denominada ordem total quando quaisquer
  dois elementos de (A) são comparáveis.

  Isto é,

  [
    (\forall a,b\in A)
    (aRb \lor bRa).
  ]

\end{defn}

\begin{defn}[Elemento Minimal]
  \label{def:matematica:relacoes:minimal}

  Seja (B\subseteq A).

  Um elemento (m\in B) é minimal quando

  [
    (\forall x\in B)
    (xRm \Rightarrow x=m).
  ]

\end{defn}

\begin{defn}[Elemento Maximal]
  \label{def:matematica:relacoes:maximal}

  Um elemento (m\in B) é maximal quando

  [
    (\forall x\in B)
    (mRx \Rightarrow x=m).
  ]

\end{defn}

\begin{defn}[Menor Elemento]
  \label{def:matematica:relacoes:menor-elemento}

  Um elemento (m\in B) é o menor elemento de (B) quando

  [
    (\forall x\in B)
    (mRx).
  ]

\end{defn}

\begin{defn}[Maior Elemento]
  \label{def:matematica:relacoes:maior-elemento}

  Um elemento (m\in B) é o maior elemento de (B) quando

  [
    (\forall x\in B)
    (xRm).
  ]

\end{defn}

\begin{prop}
  \label{prop:matematica:relacoes:menor-implica-minimal}

  Todo menor elemento é minimal.

\end{prop}

\begin{proof}

  Imediato pela definição.

\end{proof}

\begin{prop}
  \label{prop:matematica:relacoes:maior-implica-maximal}

  Todo maior elemento é maximal.

\end{prop}

\begin{proof}

  Imediato pela definição.

\end{proof}

\begin{prop}
  \label{prop:matematica:relacoes:unicidade-menor}

  Se existir um menor elemento, então ele é único.

\end{prop}

\begin{proof}

  Sejam (a) e (b) menores elementos.

  Então

  [
    aRb
    \quad\text{e}\quad
    bRa.
  ]

  Pela antissimetria,

  [
    a=b.
  ]

\end{proof}

\end{document}
